\scanpage{623}
\mspage{11}
\diagram{\begin{tikzcd}[ampersand replacement=\&,column sep=large]
\text{(1.18)} \& \left\{\vphantom{\dfrac{M}{M}}\right. M(I) \arrow[r,shift left=3,"u_!"] \arrow[r,shift right=3,"u_*","{\deleted{\unsure{u_*}}}"'] \& M(I')\arrow[l,"u^*"' ] ,
\end{tikzcd}}

on aura de façon $\pm$ tautologique

\[
\left\{\begin{aligned}
&u_!\big(\Phi\langle I\rangle\big)\subset\Phi\langle I'\rangle\\
&u^*\big(\Psi(I')\big)\subset\Psi(I).
\end{aligned}\right.\tag*{(1.19)}
\]

Le corollaire 1.8. donne : \deleted{\unsure{donc}}

\uline{Proposition} 1.10 : Supposons que \deleted{$\Phi,\Psi$} $\Phi \mathrel{\longleftarrow\!\mid\!\longrightarrow} \Psi$
soit un couple de Quillen \uline{spécial} dans $M$,
\unsure{i.e.} de la forme $(\Phi,\Phi_*)$, avec $\Phi$ partie
$Q$-close à gauche spéciale dans $M$ ($M$
satisfaisant la condition (1.10. a)). Alors
pour \unsure{tt} $I$ dans $\mathrm{Cat}$, $(\Phi\langle I\rangle,\Psi(I))$
est un couple spécial dans $M(I)$, en
particulier il satisfait au th. de factorisation.

Considérons la construction duale,
qui nous fait in\deleted{\unsure{tro}}troduire, à partir d'un
couple $\Phi\longleftrightarrow\Psi$ dans $M$, \unsure{des|les}
couples

\[
\tikz[baseline=(rel.base)]{
  \node[inner xsep=.45em,inner ysep=.6ex] (rel)
    {$\Phi(I\unsure{\rangle|)}\longleftrightarrow\Psi[I]$};
  \draw ([xshift=-1.3em,yshift=-.3ex]rel.south west)
    -- ([xshift=.25em]rel.south west) -- ([xshift=-.23em]rel.north west)
    -- ([yshift=.45ex]rel.north east) -- (rel.south east)
    -- ([xshift=.5em]rel.south east);
},\tag*{(1.20)}
\]

où $\Phi(I)$ est défini par la formule générale
(1.11) (où $\Psi$ est remplacé par $\Phi$), et où
$\Psi[I]$ est défini soit comme $\Phi(I)_*$,
soit comme

\[
\Psi[I]=\unsure{v|u}_*\big(\Psi(I_0)\big)^{-}=\Phi(I)_*.\tag*{(1.2\unsure{1|4})}
\]
% archive sheet number: 622
